% ============ 第二章 §3 伽罗瓦扩域中的素理想分解 (原书 pp.71-88) ============
\section{伽罗瓦扩域中的素理想分解}\label{sec:ch2-s3}

设数域扩张 $L/K$ 是伽罗瓦扩张.Hilbert 于 19 世纪末研究了这种扩张中素理想分解的更加精细的结构.我们在这节中介绍 Hilbert 理论的基本结果.作为应用,我们还给出分圆域中素理想分解的完整结果.

\subsection{\texorpdfstring{$n=efg$}{n=efg}}\label{sec:ch2-s3-1}

我们在第 1.3 节中定义了数域 $K$ 中理想 $\mathfrak{a}$ 的范 $N_{K/\Q}(\mathfrak{a})=|O_K/\mathfrak{a}|$.现在我们将它推广到更一般的情形(今后 $K$ 中的理想均指的是 $O_K$ 中的理想).

\emph{定义.} 设 $L/K$ 是数域的扩张,$\mathfrak{P}$ 和 $\mathfrak{p}$ 分别是 $L$ 和 $K$ 中的素理想,$\mathfrak{P}\mid\mathfrak{p}$,$f=f(\mathfrak{P}/\mathfrak{p})$,定义 $N_{L/K}(\mathfrak{P})=\mathfrak{p}^f$.更一般地,对于 $L$ 中每个理想 $\mathfrak{a}=\mathfrak{P}_1^{e_1}\cdots\mathfrak{P}_r^{e_r}$($e_i\in\Z$),定义 $N_{L/K}(\mathfrak{a})=\prod_{i=1}^{r}N_{L/K}(\mathfrak{P}_i)^{e_i}$,不难证明如此定义的(相对)范有以下简单性质(作为习题 1):

\emph{(I)} 如果 $\mathfrak{A}$ 和 $\mathfrak{B}$ 均是 $L$ 中的理想,则 $N_{L/K}(\mathfrak{AB})=N_{L/K}(\mathfrak{A})N_{L/K}(\mathfrak{B})$;

\emph{(II)} 如果 $\mathfrak{a}$ 是 $K$ 中的理想,则 $N_{L/K}(\mathfrak{a}O_L)=\mathfrak{a}^{[L:K]}$;

\emph{(III)} 在扩张 $K/\Q$ 的情形,对于 $K$ 中理想 $\mathfrak{a}$,这里关于 $N_{K/\Q}(\mathfrak{a})$ 的定义与早先的定义 $N_{K/\Q}(\mathfrak{a})=|O_K/\mathfrak{a}|$ 是一致的(但需把数 $m=|O_K/\mathfrak{a}|\in\Z$ 等同于 $\Z$ 中的理想 $(m)=m\Z$);

\emph{(IV)} 设 $L/M$ 和 $M/K$ 均是数域的扩张,则对于 $L$ 中的理想 $\mathfrak{A}$,$N_{L/K}(\mathfrak{A})=N_{M/K}(N_{L/M}(\mathfrak{A}))$.

现在设 $L/K$ 是数域的伽罗瓦扩张,$G=\Gal(L/K)$ 是该扩张的伽罗瓦群.对于每个整数 $\alpha\in O_L$,$\alpha$ 的每个 $K$-共轭元素 $\sigma(\alpha)$($\sigma\in G$)与 $\alpha$ 在 $K$ 上具有同一个极小多项式,从而均是 $L$ 中的整数,即 $\sigma(\alpha)\in O_L$.于是 $\sigma(O_L)\subseteq O_L$(对每个 $\sigma\in G$),从而必然 $\sigma(O_L)=O_L$($\sigma\in G$).类似地,若 $\mathfrak{B}$ 是 $O_L$ 的素理想,则 $\sigma(\mathfrak{B})=\{\sigma(\alpha)\mid\alpha\in\mathfrak{B}\}$ 也是 $O_L$ 的素理想,称作是 $\mathfrak{B}$ 的 $K$-共轭理想.

\begin{theorem}\label{thm:2.12}
设 $L/K$ 是数域的伽罗瓦扩张,$n=[L:K]$,$G=\Gal(L/K)$.$\mathfrak{p}$ 为 $O_K$ 中的素理想,$\mathfrak{p}$ 在 $O_L$ 中的分解式为
\[
  \mathfrak{p}O_L=\mathfrak{P}_1^{e_1}\cdots\mathfrak{P}_g^{e_g},\qquad f_i=f(\mathfrak{P}_i/\mathfrak{p})\ (1\leqslant i\leqslant g),
\]
则
\begin{enumerate}[(a)]
  \item 群 $G$ 在集合 $\{\mathfrak{P}_1,\cdots,\mathfrak{P}_g\}$ 上是可迁的,即对于每对 $\mathfrak{P}_i$ 和 $\mathfrak{P}_j$ 均存在 $\sigma\in G$,使得 $\sigma(\mathfrak{P}_i)=\mathfrak{P}_j$;
  \item 对于每个 $i$,$\{\mathfrak{P}_1,\cdots,\mathfrak{P}_g\}$ 就是 $\mathfrak{P}_i$ 的全部 $K$-共轭理想;
  \item $e_1=e_2=\cdots=e_g$,$f_1=f_2=\cdots=f_g$.令 $e=e_i$,$f=f_i$($1\leqslant i\leqslant g$),则 $n=efg$;
  \item 对于 $L$ 中每个理想 $\mathfrak{A}$,均有 $N_{L/K}(\mathfrak{A})O_L=\prod_{\sigma\in G}\sigma(\mathfrak{A})$.
\end{enumerate}
\end{theorem}

\begin{proof}
(a) 由于 $G$ 是群,我们只需证明:对于每个 $i$($1\leqslant i\leqslant g$),均存在 $\sigma\in G$,使得 $\mathfrak{P}_i=\sigma(\mathfrak{P}_1)$.证明用反证法.如果 $\mathfrak{P}_i\notin\{\sigma(\mathfrak{P}_1)\mid\sigma\in G\}$,则由中国剩余定理可知存在 $\alpha\in O_L$,使得 $\alpha\in\mathfrak{P}_i$ 并且 $\alpha\notin\sigma(\mathfrak{P}_1)$(对每个 $\sigma\in G$),于是 $\sigma(\alpha)\notin\mathfrak{P}_1$(对每个 $\sigma\in G$).因此 $N_{L/K}(\alpha)=\prod_{\sigma\in G}\sigma(\alpha)\notin\mathfrak{P}_1$,但是 $N_{L/K}(\alpha)=\prod_{\sigma\in G}\sigma(\alpha)\in\mathfrak{P}_i\cap O_K=\mathfrak{p}\subseteq\mathfrak{P}_1$,这就导致矛盾,从而证明了 (a).

(b) 我们只需证明 $\mathfrak{P}_1$ 的 $K$-共轭理想必然是 $\mathfrak{P}_1,\cdots,\mathfrak{P}_g$ 当中的一个.设 $\mathfrak{P}$ 是 $\mathfrak{P}_1$ 的一个 $K$-共轭理想,则存在 $\sigma\in G$,使得 $\mathfrak{P}=\sigma(\mathfrak{P}_1)$,但是 $\mathfrak{P}_1\cap O_K=\mathfrak{p}$,从而
\[
  \mathfrak{P}\cap O_K=\sigma(\mathfrak{P}_1)\cap\sigma(O_K)=\sigma(\mathfrak{p})=\mathfrak{p}.
\]
于是 $\mathfrak{P}\mid\mathfrak{p}$,即 $\mathfrak{P}$ 必为 $\mathfrak{P}_1,\cdots,\mathfrak{P}_g$ 中的一个.

(c) 对于每个 $i$($1\leqslant i\leqslant g$),由 (b) 知存在 $\sigma\in G$,使得 $\mathfrak{P}_i=\sigma(\mathfrak{P}_1)$.于是
\begin{align*}
  \mathfrak{P}_1^{e_1}\cdots\mathfrak{P}_g^{e_g}
  &=\mathfrak{p}O_L=\sigma(\mathfrak{p}O_L)=\sigma(\mathfrak{P}_1)^{e_1}\cdots\sigma(\mathfrak{P}_g)^{e_g}\\
  &=\mathfrak{P}_i^{e_1}\sigma(\mathfrak{P}_2)^{e_2}\cdots\sigma(\mathfrak{P}_g)^{e_g}.
\end{align*}
当 $i\geqslant 2$ 时,$\sigma(\mathfrak{P}_i)\neq\sigma(\mathfrak{P}_1)=\mathfrak{P}_i$(因为 $\mathfrak{P}_i\neq\mathfrak{P}_1$),因此由 $\mathfrak{p}O_L$ 的素理想分解式的惟一性,比较上式两边 $\mathfrak{P}_i$ 的指数,即知 $e_i=e_1$,即 $e_1=e_2=\cdots=e_g$.同样地,
\begin{align*}
  f_i&=[O_L/\mathfrak{P}_i:O_K/\mathfrak{p}]
  =[\sigma(O_L)/\sigma(\mathfrak{P}_i):\sigma(O_K)/\sigma(\mathfrak{p})]\\
  &=[O_L/\mathfrak{P}_1:O_K/\mathfrak{p}]=f_1.
\end{align*}
于是 $f_1=f_2=\cdots=f_g$.令 $e=e_i$,$f=f_i$($1\leqslant i\leqslant g$),则
\[
  n=\sum_{i=1}^{g}e_if_i=\sum_{i=1}^{g}ef=efg.
\]
并且 $\mathfrak{p}O_L=(\mathfrak{P}_1\cdots\mathfrak{P}_g)^e$.

(d) 由于范 $N_{L/K}$ 是积性函数,从而只需考虑 $L$ 中的素理想 $\mathfrak{P}$ 即可.设 $\mathfrak{P}\cap O_K=\mathfrak{p}$,$\mathfrak{p}O_L=(\mathfrak{P}_1\cdots\mathfrak{P}_g)^e$,$\mathfrak{P}=\mathfrak{P}_1$,$n=efg$.令 $D_{\mathfrak{P}}$ 为 $G$ 中固定 $\mathfrak{P}$ 的子群,即 $D_{\mathfrak{P}}=\{\sigma\in G\mid\sigma(\mathfrak{P})=\mathfrak{P}\}$.由于 $G$ 在集合 $\{\mathfrak{P}=\mathfrak{P}_1,\mathfrak{P}_2,\cdots,\mathfrak{P}_g\}$ 上是可迁的,从而 $|D_{\mathfrak{P}}|=|G|/g=n/g=ef$.并且有陪集分解
\[
  G=\sigma_1D_{\mathfrak{P}}\cup\sigma_2D_{\mathfrak{P}}\cup\cdots\cup\sigma_gD_{\mathfrak{P}},\qquad
  \sigma_i(\mathfrak{P})=\mathfrak{P}_i\ (1\leqslant i\leqslant g).
\]
于是
\begin{align*}
  \prod_{\sigma\in G}\sigma(\mathfrak{P})
  &=\prod_{i=1}^{g}\sigma_i\Big(\prod_{\sigma\in D_{\mathfrak{P}}}\sigma(\mathfrak{P})\Big)
  =\prod_{i=1}^{g}\sigma_i(\mathfrak{P}^{ef})\\
  &=(\mathfrak{P}_1\cdots\mathfrak{P}_g)^{ef}=(\mathfrak{p}O_L)^f=N_{L/K}(\mathfrak{P})O_L.
\end{align*}
\end{proof}

\subsection{分解群和惯性群}\label{sec:ch2-s3-2}

我们在定理 \ref{thm:2.12} 的证明中曾经指出,对于数域的伽罗瓦扩张 $L/K$ 和域 $L$ 的每个素理想 $\mathfrak{P}$,
\[
  D_{\mathfrak{P}}=\{\sigma\in\Gal(L/K)\mid\sigma(\mathfrak{P})=\mathfrak{P}\}
\]
是 $\Gal(L/K)$ 的 $ef$ 阶子群,称作是 $\mathfrak{P}$(对于伽罗瓦扩张 $L/K$)的\emph{分解群}.令 $\overline{L}=O_L/\mathfrak{P}$,$\overline{K}=O_K/\mathfrak{p}$(其中 $\mathfrak{p}=\mathfrak{P}\cap O_K$),则 $\overline{L}$ 和 $\overline{K}$ 均是有限域,并且 $\overline{L}$ 是 $\overline{K}$ 的 $f$ 次扩域.对于每个 $\sigma\in D_{\mathfrak{P}}$,由于 $\sigma(\mathfrak{P})=\mathfrak{P}$,从而可以定义如下的映射:
\[
  \overline{\sigma}\colon\overline{L}\to\overline{L},\qquad
  \overline{\sigma}(\overline{a})=\overline{\sigma(a)}\quad (a\in O_L,\ \overline{a}=a+\mathfrak{P}).
\]
这个映射有如下的性质:

(1) 首先,映射 $\overline{\sigma}$ 是确切定义的.因为若 $a,b\in O_L$,$\overline{a}=\overline{b}$,则 $a-b\in\mathfrak{P}$,从而 $\sigma(a)-\sigma(b)=\sigma(a-b)\in\sigma(\mathfrak{P})=\mathfrak{P}$,即 $\overline{\sigma(a)}=\overline{\sigma(b)}$,即 $\overline{\sigma(a)}$ 与剩余类 $\overline{a}$ 中代表元 $a$ 的选取无关.

(2) $\overline{\sigma}$ 是域 $\overline{L}$ 的自同构.因为
\begin{align*}
  \overline{\sigma}(\overline{a}\pm\overline{b})&=\overline{\sigma(a\pm b)}
  =\overline{\sigma(a)\pm\sigma(b)}=\overline{\sigma(a)}\pm\overline{\sigma(b)}\\
  &=\overline{\sigma}(\overline{a})\pm\overline{\sigma}(\overline{b}),
\end{align*}
类似地,$\overline{\sigma}(\overline{a}\cdot\overline{b})=\overline{\sigma}(\overline{a})\cdot\overline{\sigma}(\overline{b})$.从而 $\overline{\sigma}$ 是域 $\overline{L}$ 的自同态.进而
\begin{align*}
  \operatorname{Ker}\overline{\sigma}&=\{\overline{a}\mid\overline{\sigma(a)}=\overline{0}\}
  =\{\overline{a}\mid\sigma(a)\in\mathfrak{P}\}\\
  &=\{\overline{a}\mid a\in\sigma^{-1}(\mathfrak{P})=\mathfrak{P}\}=\{\overline{0}\},
\end{align*}
从而 $\overline{\sigma}$ 是单同态.最后,由于 $\overline{L}$ 是有限域,从而 $\overline{\sigma}$ 也必然是满同态,即 $\overline{\sigma}$ 是域 $\overline{L}$ 的自同构.

(3) $\overline{\sigma}$ 固定 $\overline{K}$ 中每个元素.这是因为:若 $\overline{a}\in\overline{K}$,$a\in O_K$,则 $\overline{\sigma}(\overline{a})=\overline{\sigma(a)}=\overline{a}$.

由以上即知,对于每个 $\sigma\in D_{\mathfrak{P}}$,我们按上述方式定义出的 $\overline{\sigma}$ 属于 $\overline{L}/\overline{K}$ 的伽罗瓦群 $\Gal(\overline{L}/\overline{K})$.从而又给出映射:
\[
  \pi\colon D_{\mathfrak{P}}\to\Gal(\overline{L}/\overline{K}),\qquad \pi(\sigma)=\overline{\sigma}.
\]

\begin{lemma}\label{lem:ch2-15}
$\pi$ 是群的满同态,并且 $\operatorname{Ker}\pi$ 为
\[
  I_{\mathfrak{P}}=\{\sigma\in D_{\mathfrak{P}}\mid\sigma(\alpha)\equiv\alpha\ (\mathrm{mod}\ \mathfrak{P})\text{,对每个 }\alpha\in O_L\}.
\]
$|I_{\mathfrak{P}}|=e$,$D_{\mathfrak{P}}/I_{\mathfrak{P}}$ 为 $f$ 阶循环群.
\end{lemma}

\begin{proof}
设 $\sigma,\tau\in D_{\mathfrak{P}}$,则对每个 $\overline{\alpha}\in\overline{L}$,$\alpha\in O_L$,均有
\[
  \overline{\sigma\tau}(\bar{\alpha})=\overline{\sigma\tau(\alpha)}
  =\overline{\sigma}(\overline{\tau(\alpha)})=\overline{\sigma}\,\overline{\tau}(\bar{\alpha}),
\]
因此 $\overline{\sigma\tau}=\overline{\sigma}\,\overline{\tau}$,这表明 $\pi$ 是群的同态.为证 $\pi$ 是满同态,只需证明对每个 $\widetilde{\sigma}\in\Gal(\overline{L}/\overline{K})$,均存在 $\sigma\in D_{\mathfrak{P}}$,使得 $\overline{\sigma}=\widetilde{\sigma}$.我们知道,有限域扩张 $\overline{L}/\overline{K}$ 均是单扩张 $\overline{L}=\overline{K}(\overline{\alpha})$,$\alpha\in O_L$(附录 A, IV).以 $K_D$ 表示 $L$ 的子域,它在伽罗瓦对应下对应于群 $D_{\mathfrak{P}}$,即
\[
  K_D=\{\alpha\in L\mid\sigma(\alpha)=\alpha\text{,对于每个 }\sigma\in D_{\mathfrak{P}}\}.
\]
从伽罗瓦扩张基本定理(附录 A, (16))知道 $L\supseteq K_D\supseteq K$,$L/K_D$ 是 $ef$ 次伽罗瓦扩张,并且 $\Gal(L/K_D)=D_{\mathfrak{P}}$.令 $O_D=O_{K_D}$,$\mathfrak{P}\cap O_D=\mathfrak{P}_D$,则 $\mathfrak{P}_D\cap O_K=\mathfrak{P}\cap O_K=\mathfrak{p}$.从而我们有下边的图表:
\[
  \begin{array}{ccc}
    \mathfrak{P} & L & \{1\}\\
    \mid & \mid & \mid\\
    \mathfrak{P}_D & K_D & D_{\mathfrak{P}}\\
    \mid & \mid & \mid\\
    \mathfrak{p} & K & \Gal(L/K)
  \end{array}\Bigg\} g
\]
由于 $\Gal(L/K_D)=D_{\mathfrak{P}}$,而对 $\sigma\in D_{\mathfrak{P}}$ 均有 $\sigma(\mathfrak{P})=\mathfrak{P}$.从定理 \ref{thm:2.12} (b) 即知 $\mathfrak{P}$ 对于扩张 $L/K_D$ 是自共轭的.因此 $\mathfrak{P}_D$ 在 $L$ 中的素理想分解为 $\mathfrak{P}_DO_L=\mathfrak{P}^{e'}$,$e'=e(\mathfrak{P}/\mathfrak{P}_D)$.令 $f'=f(\mathfrak{P}\mid\mathfrak{P}_D)$,则
\[
  e'\cdot e(\mathfrak{P}_D\mid\mathfrak{p})=e(\mathfrak{P}/\mathfrak{p})=e,\qquad
  f'\cdot f(\mathfrak{P}_D\mid\mathfrak{p})=f(\mathfrak{P}/\mathfrak{p})=f.
\]
但是 $e'f'=[L:K_D]=ef$,这就表明 $e'=e$,$f'=f$,并且 $e(\mathfrak{P}_D/\mathfrak{p})=f(\mathfrak{P}_D/\mathfrak{p})=1$.特别地我们得到 $O_D/\mathfrak{P}_D=O_K/\mathfrak{p}$.令 $\alpha$ 在 $K_D$ 上的极小多项式为
\[
  f(x)=x^r+a_1x^{r-1}+\cdots+a_r,\qquad a_i\in O_D,
\]
则 $f(x)$ 的每个根均有形式 $\sigma(\alpha)$,$\sigma\in\Gal(L/K_D)=D_{\mathfrak{P}}$.将 $f(x)$ 的系数模 $\mathfrak{P}_D$ 之后,得到 $\overline{f}(x)=x^r+\overline{a}_1x^{r-1}+\cdots+\overline{a}_r\in O_D/\mathfrak{P}_D[x]$.由于 $O_D/\mathfrak{P}_D=O_K/\mathfrak{p}=\overline{K}$,从而 $\overline{a}_i$ 中可取代表元(仍记为 $a_i$)属于 $O_K$($1\leqslant i\leqslant r$).因此可认为 $\overline{f}(x)\in\overline{K}[x]$.而这时 $\overline{f}(x)$ 的每个根均有形式 $\overline{\sigma(\alpha)}=\overline{\sigma}(\overline{\alpha})$.由于 $\overline{\alpha}$ 在 $\overline{K}$ 上的极小多项式为 $\overline{f}(x)$ 的因式,从而 $\overline{\alpha}$ 的 $\overline{K}$-共轭元 $\widetilde{\sigma}(\overline{\alpha})$ 是 $\overline{f}(x)$ 的根.从而 $\widetilde{\sigma}(\overline{\alpha})$ 必有形式 $\overline{\sigma(\alpha)}=\overline{\sigma}(\overline{\alpha})$.这就表明存在 $\sigma\in D_{\mathfrak{P}}$,使得 $\overline{\sigma}(\overline{\alpha})=\widetilde{\sigma}(\overline{\alpha})$.由于 $\overline{L}=\overline{K}(\overline{\alpha})$,因此 $\overline{\sigma}=\widetilde{\sigma}$.即 $\pi\colon D_{\mathfrak{P}}\to\Gal(\overline{L}/\overline{K})$ 是满同态.

最后,
\begin{align*}
  \operatorname{Ker}\pi&=\{\sigma\in D_{\mathfrak{P}}\mid\overline{\sigma}=1\}\\
  &=\{\sigma\in D_{\mathfrak{P}}\mid\overline{\sigma}(\overline{\alpha})=\overline{\alpha}\text{,对于每个 }\overline{\alpha}\in\overline{L}\}\\
  &=\{\sigma\in D_{\mathfrak{P}}\mid\sigma(\alpha)\equiv\alpha\ (\mathrm{mod}\ \mathfrak{P})\text{,对于每个 }\alpha\in O_L\}=I_{\mathfrak{P}}.
\end{align*}
从而 $D_{\mathfrak{P}}/I_{\mathfrak{P}}\cong\Gal(\overline{L}/\overline{K})$.但是右边为 $f$ 阶循环群(附录 A, (17)),从而 $D_{\mathfrak{P}}/I_{\mathfrak{P}}$ 也是如此,这就完成了引理的证明.
\end{proof}

引理 \ref{lem:ch2-15} 中的群 $I_{\mathfrak{P}}$ 叫作是 $\mathfrak{P}$(对于伽罗瓦扩张 $L/K$)的\emph{惯性群}.这是分解群 $D_{\mathfrak{P}}$ 的 $e$ 阶正规子群,并且 $D_{\mathfrak{P}}/I_{\mathfrak{P}}$ 为 $f$ 阶循环群.在伽罗瓦对应下,分解群 $D_{\mathfrak{P}}$ 对应的域
\[
  K_D=\{\alpha\in L\mid\sigma(\alpha)=\alpha\text{,对每个 }\sigma\in D_{\mathfrak{P}}\}
\]
叫作是 $\mathfrak{P}$ 的\emph{分解域},而惯性群 $I_{\mathfrak{P}}$ 对应的域
\[
  K_I=\{\alpha\in L\mid\sigma(\alpha)=\alpha\text{,对每个 }\sigma\in I_{\mathfrak{P}}\}
\]
叫作是 $\mathfrak{P}$ 的\emph{惯性域}.于是我们有 $K\subseteq K_D\subseteq K_I\subseteq L$,$[L:K_I]=e$,$[K_I:K_D]=|D_{\mathfrak{P}}/I_{\mathfrak{P}}|=f$,$[K_D:K]=g$.

\begin{theorem}\label{thm:2.13}
设 $L/K$ 为数域的伽罗瓦扩张,$n=[L:K]$.$K$ 中素理想 $\mathfrak{p}$ 在 $L$ 中的素理想分解式为
\[
  \mathfrak{p}O_L=(\mathfrak{P}_1\cdots\mathfrak{P}_g)^e,\qquad f=f(\mathfrak{P}_i/\mathfrak{p})\ (1\leqslant i\leqslant g),\qquad n=efg,\quad \mathfrak{P}=\mathfrak{P}_1.
\]
令 $D_{\mathfrak{P}}$,$I_{\mathfrak{P}}$,$K_D$,$K_I$ 分别是 $\mathfrak{P}$ 对于伽罗瓦扩张 $L/K$ 的分解群,惯性群,分解域和惯性域,则我们有如下的图表:
\[
  \begin{array}{cccc}
    (\text{素理想分解}) & (\text{域}) & (\text{群}) & (\text{剩余类域})\\
    \mathfrak{P} & L & \{1\}\Bigg\}_e & \overline{L}=O_L/\mathfrak{P}\\
    \mid & \mid & \mid & \parallel\\
    \mathfrak{P}_I & K_I & I_{\mathfrak{P}} & \\
    \mid & \mid & \mid & \\
    \mathfrak{P}_D & K_D & D_{\mathfrak{P}} & \\
    \mid & \mid & \mid & \\
    \mathfrak{p} & K & \Gal(L/K) & \overline{K}=O_K/\mathfrak{p}
  \end{array}
\]
具体说来就是:

(a) $L/K_I$,$L/K_D$ 和 $K_I/K_D$ 均是伽罗瓦扩张,其伽罗瓦群分别为 $I_{\mathfrak{P}}$,$D_{\mathfrak{P}}$ 和 $D_{\mathfrak{P}}/I_{\mathfrak{P}}\cong\Gal(\overline{L}/\overline{K})$.

(b) 令 $\mathfrak{P}_I=\mathfrak{P}\cap O_I$,$\mathfrak{P}_D=\mathfrak{P}\cap O_D$,则:

(i) $\mathfrak{p}=\mathfrak{P}_D\cap O_K$,$e(\mathfrak{P}_D/\mathfrak{p})=f(\mathfrak{P}_D/\mathfrak{p})=1$,并且当 $D_{\mathfrak{P}}$ 为 $\Gal(L/K)$ 的正规子群时,$\mathfrak{p}$ 在 $K_D$ 中完全分裂;

(ii) $\mathfrak{P}_DO_I=\mathfrak{P}_I$,$f(\mathfrak{P}_I/\mathfrak{P}_D)=f$(即 $\mathfrak{P}_D$ 在 $K_I$ 中是惯性的,并且剩余类域次数为 $f$);

(iii) $\mathfrak{P}_IO_L=\mathfrak{P}^e$,$f(\mathfrak{P}/\mathfrak{P}_I)=1$(即 $\mathfrak{P}_I$ 在 $L$ 中完全分歧).

(c) 对于每个 $\sigma\in\Gal(L/K)$,
\[
  D_{\sigma(\mathfrak{P})}=\sigma D_{\mathfrak{P}}\sigma^{-1},\qquad
  I_{\sigma(\mathfrak{P})}=\sigma I_{\mathfrak{P}}\sigma^{-1},\qquad
  K_{I\sigma(\mathfrak{P})}=\sigma(K_{I\mathfrak{P}}),\qquad
  K_{D\sigma(\mathfrak{P})}=\sigma(K_{D\mathfrak{P}}).
\]
\end{theorem}

\begin{proof}
(a) 由伽罗瓦理论即可证得(注意:$K_I/K_D$ 为伽罗瓦扩张是由于 $I_{\mathfrak{P}}$ 为 $D_{\mathfrak{P}}$ 的正规子群).

(b) 由于
\[
  e=e(\mathfrak{P}/\mathfrak{p})=e(\mathfrak{P}/\mathfrak{P}_I)e(\mathfrak{P}_I/\mathfrak{P}_D)e(\mathfrak{P}_D/\mathfrak{p}),
\]
\[
  f=f(\mathfrak{P}/\mathfrak{p})=f(\mathfrak{P}/\mathfrak{P}_I)\cdot f(\mathfrak{P}_I/\mathfrak{P}_D)\cdot f(\mathfrak{P}_D/\mathfrak{p}).
\]
我们在引理 \ref{lem:ch2-15} 的证明中已经得出 $e(\mathfrak{P}_D/\mathfrak{p})=f(\mathfrak{P}_D/\mathfrak{p})=1$,即
\[
  e(\mathfrak{P}/\mathfrak{P}_D)=e(\mathfrak{P}/\mathfrak{P}_I)\cdot e(\mathfrak{P}_I/\mathfrak{P}_D)=e,\qquad
  f(\mathfrak{P}/\mathfrak{P}_I)f(\mathfrak{P}_I/\mathfrak{P}_D)=f.
\]
并且当 $D_{\mathfrak{P}}$ 为 $\Gal(L/K)$ 的正规子群时,$K_D/K$ 为伽罗瓦扩张,即知 $\mathfrak{p}$ 在 $K_D$ 中完全分裂.继而考虑伽罗瓦扩张 $L/K_I$.由定义不难看出 $\mathfrak{P}$ 对于扩张 $L/K_I$ 的分解群和惯性群均是 $I_{\mathfrak{P}}$,因此 $f(\mathfrak{P}/\mathfrak{P}_I)=|I_{\mathfrak{P}}/I_{\mathfrak{P}}|=1$.于是 $f(\mathfrak{P}_I/\mathfrak{P}_D)=f$,并且由 $e=e(\mathfrak{P}/\mathfrak{P}_I)\cdot e(\mathfrak{P}_I/\mathfrak{P}_D)$ 及 $e(\mathfrak{P}/\mathfrak{P}_I)=[L:K_I]/f(\mathfrak{P}/\mathfrak{P}_I)=e/1=e$ 可知 $e(\mathfrak{P}_I/\mathfrak{P}_D)=1$.

(c) 由于
\[
  \tau\in D_{\mathfrak{P}}\Leftrightarrow\tau(\mathfrak{P})=\mathfrak{P}
  \Leftrightarrow\sigma\tau\sigma^{-1}(\sigma(\mathfrak{P}))=\sigma(\mathfrak{P})
  \Leftrightarrow\sigma\tau\sigma^{-1}\in D_{\sigma(\mathfrak{P})},
\]
这就表明 $D_{\sigma(\mathfrak{P})}=\sigma D_{\mathfrak{P}}\sigma^{-1}$.类似可证 $I_{\sigma(\mathfrak{P})}=\sigma I_{\mathfrak{P}}\sigma^{-1}$.另一方面,
\begin{align*}
  \alpha\in K_{D\mathfrak{P}}
  &\Longleftrightarrow\tau(\alpha)=\alpha\ (\text{对每个 }\tau\in D_{\mathfrak{P}})\\
  &\Longleftrightarrow\sigma\tau\sigma^{-1}(\sigma(\alpha))=\sigma(\alpha)\ (\text{对每个 }\sigma\tau\sigma^{-1}\in\sigma D_{\mathfrak{P}}\sigma^{-1}=D_{\sigma(\mathfrak{P})})\\
  &\Longleftrightarrow\sigma(\alpha)\in K_{D\sigma(\mathfrak{P})},
\end{align*}
从而 $K_{D\sigma(\mathfrak{P})}=\sigma(K_{D\mathfrak{P}})$.同样可证 $K_{I\sigma(\mathfrak{P})}=\sigma(K_{I\mathfrak{P}})$.
\end{proof}

\begin{remark}
设 $L/K$ 为数域的伽罗瓦扩张.如果其伽罗瓦群 $\Gal(L/K)$ 为 Abel 群,我们也称 $L/K$ 为 Abel 扩张.类似地,如果 $\Gal(L/K)$ 是循环群,也称 $L/K$ 是循环扩张.从定理 \ref{thm:2.13} 可以看出,对于 Abel 扩张 $L/K$,若 $K$ 的素理想 $\mathfrak{p}$ 在 $L$ 中分解为 $\mathfrak{p}=(\mathfrak{P}_1\cdots\mathfrak{P}_g)^e$,则 $D_{\mathfrak{P}_i}$($1\leqslant i\leqslant g$)均相等.我们可以将它表示成 $D_{\mathfrak{p}}$,叫作是 $\mathfrak{p}$(对于 $L/K$)的分解群.类似地,$I_{\mathfrak{P}_i}$($1\leqslant i\leqslant g$)也均相等,表示成 $I_{\mathfrak{p}}$,叫作是 $\mathfrak{p}$(对于 $L/K$)的惯性群.在伽罗瓦对应下,$D_{\mathfrak{p}}$ 和 $I_{\mathfrak{p}}$ 对应的域分别为 $K_D$ 和 $K_I$,后者也分别称作是 $\mathfrak{p}$ 的分解域和惯性域.于是当 $L/K$ 为 Abel 扩张时,可以把定理 \ref{thm:2.13} 粗略地叙述为:$\mathfrak{p}$ 在分解域 $K_D$ 中完全分裂:$\mathfrak{p}O_D=\mathfrak{P}_{D,1}\cdots\mathfrak{P}_{D,g}$;每个 $\mathfrak{P}_{D,i}$ 在惯性域 $K_I$ 中均惯性:$\mathfrak{P}_{D,i}O_I=\mathfrak{P}_{I,i}$($1\leqslant i\leqslant g$),$f(\mathfrak{P}_{I,i}/\mathfrak{P}_{D,i})=f$;最后每个 $\mathfrak{P}_{I,i}$ 在 $L$ 中均完全分歧:$\mathfrak{P}_{I,i}O_L=\mathfrak{P}_i^e$,总的结果为 $\mathfrak{p}O_L=(\mathfrak{P}_1\cdots\mathfrak{P}_g)^e$.
\end{remark}

\subsection{Frobenius 自同构}\label{sec:ch2-s3-3}

设 $L/K$ 为数域的伽罗瓦扩张,如果 $K$ 的素理想 $\mathfrak{p}$ 在 $L$ 中不分歧,即 $\mathfrak{p}O_L=\mathfrak{P}_1\cdots\mathfrak{P}_g$,$\mathfrak{P}=\mathfrak{P}_1$,则由定理 \ref{thm:2.13} 可知 $I_{\mathfrak{P}}=\{1\}$,而 $D_{\mathfrak{P}}$ 是 $f$ 阶循环群,并且 $D_{\mathfrak{P}}$ 自然同构于剩余类域的伽罗瓦群 $\Gal(\overline{L}/\overline{K})$.注意 $|\overline{K}|=|O_K/\mathfrak{p}|=N_{K/\Q}(\mathfrak{p})$(以下简记作 $N(\mathfrak{p})$).从而 $\Gal(\overline{L}/\overline{K})$ 是由 $f$ 阶元素 $\sigma\colon\gamma\mapsto\gamma^{N(\mathfrak{p})}$($\gamma\in\overline{L}$)所生成的(附录 A, (17)).在自然同构之下,$D_{\mathfrak{P}}$ 中对应于 $\sigma$ 的元素表示成 $\left(\dfrac{L/K}{\mathfrak{P}}\right)$,它是 $f$ 阶元素并且生成 $D_{\mathfrak{P}}$.我们将 $\left(\dfrac{L/K}{\mathfrak{P}}\right)$ 叫作是 $\mathfrak{P}$ 对于 $L/K$ 的 \emph{Frobenius 自同构}.$D_{\mathfrak{P}}$ 中的 Frobenius 自同构这个元素可以用
\[
  \Big(\frac{L/K}{\mathfrak{P}}\Big)\alpha\equiv\alpha^{N(\mathfrak{p})}\ (\mathrm{mod}\ \mathfrak{P})\quad (\text{对每个 }\alpha\in O_L)
\]
来刻画.这是因为:如果 $\tau\in D_{\mathfrak{P}}$ 并且对每个 $\alpha\in O_L$ 均有 $\tau\alpha\equiv\alpha^{N(\mathfrak{p})}\ (\mathrm{mod}\ \mathfrak{P})$,则在自然同构 $D_{\mathfrak{P}}\cong\Gal(\overline{L}/\overline{K})$,$\tau\mapsto\overline{\tau}$ 之下,对每个 $\overline{\alpha}\in\overline{L}$ 均有 $\overline{\tau}(\overline{\alpha})=\overline{\alpha}^{N(\mathfrak{p})}$,即 $\overline{\tau}=\sigma$,从而 $\tau=\left(\dfrac{L/K}{\mathfrak{P}}\right)$.

Frobenius 自同构的基本性质为:

\begin{lemma}\label{lem:ch2-16}
设 $L/K$ 为数域的伽罗瓦扩张,$\mathfrak{P}\mid\mathfrak{p}$,其中 $\mathfrak{P}$ 和 $\mathfrak{p}$ 分别为 $L$ 和 $K$ 的素理想,并且 $e(\mathfrak{P}/\mathfrak{p})=1$,则

(a) 对于每个 $\sigma\in\Gal(L/K)$,
\[
  \Big(\frac{L/K}{\sigma(\mathfrak{P})}\Big)=\sigma\Big(\frac{L/K}{\mathfrak{P}}\Big)\sigma^{-1};
\]
(b) 如果 $E$ 是 $L/K$ 的中间域(从而 $L/E$ 为伽罗瓦扩张),$\mathfrak{P}\cap O_E=\mathfrak{P}_E$,则 $e(\mathfrak{P}\mid\mathfrak{P}_E)=1$,并且
\[
  \Big(\frac{L/E}{\mathfrak{P}}\Big)=\Big(\frac{L/K}{\mathfrak{P}}\Big)^{f(\mathfrak{P}_E/\mathfrak{p})};
\]
(c) 如果 $E/K$ 也是伽罗瓦扩张,则 $e(\mathfrak{P}_E/\mathfrak{p})=1$,并且
\[
  \Big(\frac{E/K}{\mathfrak{P}_E}\Big)=\Big(\frac{L/K}{\mathfrak{P}}\Big)\Big|_E
  \quad\Big(\Big(\frac{L/K}{\mathfrak{P}}\Big)\text{在 }E\text{ 上的限制}\Big).
\]
\end{lemma}

\begin{proof}
(a)
\[
  \Big(\frac{L/K}{\mathfrak{P}}\Big)\alpha\equiv\alpha^{N(\mathfrak{p})}\ (\mathrm{mod}\ \mathfrak{P})\ (\alpha\in O_L)
\]
\[
  \Rightarrow\sigma\Big(\frac{L/K}{\mathfrak{P}}\Big)\sigma^{-1}(\sigma(\alpha))
  \equiv(\sigma(\alpha))^{N(\mathfrak{p})}\ (\mathrm{mod}\ \sigma(\mathfrak{P}))\ (\alpha\in O_L)
\]
\[
  \Rightarrow\sigma\Big(\frac{L/K}{\mathfrak{P}}\Big)\sigma^{-1}\alpha
  \equiv\alpha^{N(\mathfrak{p})}\ (\mathrm{mod}\ \sigma(\mathfrak{P}))\ (\alpha\in O_L)
\]
\[
  \Rightarrow\sigma\Big(\frac{L/K}{\mathfrak{P}}\Big)\sigma^{-1}=\Big(\frac{L/K}{\sigma(\mathfrak{P})}\Big).
\]

(b)
\begin{align*}
  N_{E/\Q}(\mathfrak{P}_E)&=N_{K/\Q}(\mathfrak{p})^{f(\mathfrak{P}_E/\mathfrak{p})}\\
  &\Rightarrow\Big(\frac{L/K}{\mathfrak{P}}\Big)^{f(\mathfrak{P}_E/\mathfrak{p})}\alpha
  \equiv\alpha^{N_{K/\Q}(\mathfrak{p})^{f(\mathfrak{P}_E/\mathfrak{p})}}
  =\alpha^{N_{E/\Q}(\mathfrak{P}_E)}\ (\mathrm{mod}\ \mathfrak{P})\ (\alpha\in O_L)\\
  &\Rightarrow\Big(\frac{L/K}{\mathfrak{P}}\Big)^{f(\mathfrak{P}_E/\mathfrak{p})}
  =\Big(\frac{L/E}{\mathfrak{P}}\Big).
\end{align*}

(c) 由于 $\left(\dfrac{L/K}{\mathfrak{P}}\right)\alpha\equiv\alpha^{N(\mathfrak{p})}\ (\mathrm{mod}\ \mathfrak{P})$($\alpha\in O_L$),而 $\mathfrak{P}\cap O_E=\mathfrak{P}_E$,从而对 $\alpha\in O_E$ 有
\[
  \Big(\frac{L/K}{\mathfrak{P}}\Big)\alpha\equiv\alpha^{N(\mathfrak{p})}\ (\mathrm{mod}\ \mathfrak{P}_E),
\]
这就表明 $\left(\dfrac{L/K}{\mathfrak{P}}\right)\Big|_E=\left(\dfrac{E/K}{\mathfrak{P}_E}\right)$.
\end{proof}

作为引理 \ref{lem:ch2-16} 的应用,我们有如下的引理 \ref{lem:ch2-17}.

\begin{lemma}\label{lem:ch2-17}
(a) 设 $E_1/K$,$E_2/K$ 均是数域的伽罗瓦扩张,$L=E_1E_2$(域的合成,见附录 A, III),则 $L/K$ 也是伽罗瓦扩张,并且 $\Gal(L/K)$ 同构于直积 $\Gal(E_1/K)\times\Gal(E_2/K)$ 的一个子群.

(b) $K$ 中素理想 $\mathfrak{p}$ 在 $L$ 中不分歧 $\Leftrightarrow\mathfrak{p}$ 在 $E_1$ 和 $E_2$ 中均不分歧.

(c) $\mathfrak{p}$ 在 $L$ 中完全分裂 $\Leftrightarrow\mathfrak{p}$ 在 $E_1$ 和 $E_2$ 中均完全分裂.
\end{lemma}

\begin{proof}
(a) $E_i/K$($i=1,2$)为伽罗瓦扩张 $\Rightarrow E_i$ 是某个多项式 $f_i(x)\in K[x]$ 的分裂域($i=1,2$)$\Rightarrow L=E_1E_2$ 是多项式 $f_1(x)f_2(x)\in K[x]$ 的分裂域 $\Rightarrow L/K$ 是伽罗瓦扩张.进而,作映射
\[
  \varphi\colon\Gal(L/K)\to\Gal(E_1/K)\times\Gal(E_2/K),\qquad
  \sigma\mapsto(\sigma|_{E_1},\sigma|_{E_2}),
\]
易知 $\varphi$ 为群的单同态.从而 $\Gal(L/K)$ 同构于 $\Gal(E_1/K)\times\Gal(E_2/K)$ 的子群 $\varphi(\Gal(L/K))$.

(b) 设 $\mathfrak{P}\mid\mathfrak{p}$,$\mathfrak{P}$ 为 $L$ 的素理想.又令 $\mathfrak{P}_i=\mathfrak{P}\cap O_{E_i}$($i=1,2$).若 $\sigma\in D_{\mathfrak{P}}$,则 $\sigma(\mathfrak{P})=\mathfrak{P}$.从而
\[
  \sigma(\mathfrak{P}_i)=\sigma(\mathfrak{P}\cap O_{E_i})=\sigma(\mathfrak{P})\cap\sigma(O_{E_i})=\mathfrak{P}\cap O_{E_i}=\mathfrak{P}_i.
\]
因此 $\sigma|_{E_i}\in D_{\mathfrak{P}_i}$($i=1,2$),即 $\varphi(D_{\mathfrak{P}})\subseteq D_{\mathfrak{P}_1}\times D_{\mathfrak{P}_2}$.同样可证 $\varphi(I_{\mathfrak{P}})\subseteq I_{\mathfrak{P}_1}\times I_{\mathfrak{P}_2}$.从而若 $\mathfrak{p}$ 在 $E_i$ 中不分歧($i=1,2$),即 $I_{\mathfrak{P}_1}=I_{\mathfrak{P}_2}=\{1\}$,则 $I_{\mathfrak{P}}=\{1\}$,即 $\mathfrak{p}$ 在 $L$ 中不分歧.反之,由 $\mathfrak{p}$ 在 $L$ 中不分歧显然可推得 $\mathfrak{p}$ 在 $L$ 的子域 $E_i$($i=1,2$)中也不分歧(第 2 节习题 5).

(c) 从引理 \ref{lem:ch2-16} 知道,
\begin{align*}
  \varphi\Big(\Big(\frac{L/K}{\mathfrak{P}}\Big)\Big)
  &=\Big(\Big(\frac{L/K}{\mathfrak{P}}\Big)\Big|_{E_1},\Big(\frac{L/K}{\mathfrak{P}}\Big)\Big|_{E_2}\Big)\\
  &=\Big(\Big(\frac{E_1/K}{\mathfrak{P}_1}\Big),\Big(\frac{E_2/K}{\mathfrak{P}_2}\Big)\Big).
\end{align*}
从而
\begin{align*}
  \mathfrak{p}\text{ 在 }L\text{ 中完全分裂}
  &\Leftrightarrow e(\mathfrak{P}/\mathfrak{p})=f(\mathfrak{P}/\mathfrak{p})=1
  \Leftrightarrow\Big(\frac{L/K}{\mathfrak{P}}\Big)=1\\
  &\Leftrightarrow\Big(\frac{E_1/K}{\mathfrak{P}_1}\Big)=1,\ \Big(\frac{E_2/K}{\mathfrak{P}_2}\Big)=1
  \Leftrightarrow\mathfrak{p}\text{ 在 }E_1\text{ 和 }E_2\text{ 中均完全分裂}.
\end{align*}
\end{proof}

作为引理 \ref{lem:ch2-16} 的另一个应用,我们现在证明

\begin{lemma}\label{lem:ch2-18}
设 $L/K$ 是数域的 Abel 扩张,$\mathfrak{p}$ 是 $O_K$ 的素理想,$K_D$ 和 $K_I$ 分别是 $\mathfrak{p}$ 对于扩张 $L/K$ 的分解域和惯性域,则 $K_D$ 是使 $\mathfrak{p}$ 完全分裂的最大中间域,而 $K_I$ 是使 $\mathfrak{p}$ 不分歧的最大中间域.换句话说,若 $K\subseteq M\subseteq L$,如果 $\mathfrak{p}$ 在 $M$ 中完全分裂,则 $M\subseteq K_D$;如果 $\mathfrak{p}$ 在 $M$ 中不分歧,则 $M\subseteq K_I$.
\end{lemma}

\begin{proof}
由于 $L/K$ 是 Abel 扩张,从而 $K_I$,$K_D$ 和 $M$ 也是 $K$ 的伽罗瓦扩张.令 $L'=MK_D$.由定理 \ref{thm:2.13} 可知
\[
  [K_D:K]=\mathfrak{p}\text{ 在 }K_D\text{ 中的分裂次数}
  =\mathfrak{p}\text{ 在 }L\text{ 中的分裂次数}
  \geqslant\mathfrak{p}\text{ 在 }L'\text{ 中的分裂次数}.
\]
如果 $\mathfrak{p}$ 在 $M$ 中完全分裂,由引理 \ref{lem:ch2-17} 可知 $\mathfrak{p}$ 在 $L'$ 中也完全分裂,于是
\[
  \mathfrak{p}\text{ 在 }L'\text{ 中的分裂次数}=[L':K]\geqslant[K_D:K].
\]
综合上述可知 $[L':K]=[K_D:K]$.但是 $L'=MK_D\supseteq K_D$,于是 $MK_D=K_D$,即 $M\subseteq K_D$.这就表明 $K_D$ 是使 $\mathfrak{p}$ 完全分裂的最大中间域.同样可证 $K_I$ 是使 $\mathfrak{p}$ 不分歧的最大中间域.
\end{proof}

\begin{example}\label{ex:ch2-s3-3-1}
对每个素数 $p$,决定 $p$ 在 $K=\Q(\sqrt{-1},\sqrt{3})$ 中的惯性域和分解域.

\emph{解} $K/\Q$ 是四次伽罗瓦扩张,而伽罗瓦群 $G=\Gal(K/\Q)$ 是由 $\sigma$ 和 $\tau$ 生成的两个 $2$ 阶循环群的直积,其中
\[
  \sigma(\sqrt{-1})=-\sqrt{-1},\quad \sigma(\sqrt{3})=\sqrt{3},\ \text{从而 }\sigma(\sqrt{-3})=-\sqrt{-3};
\]
\[
  \tau(\sqrt{-1})=\sqrt{-1},\quad \tau(\sqrt{3})=-\sqrt{3},\ \text{从而 }\tau(\sqrt{-3})=-\sqrt{-3}.
\]
于是 $G=\{I,\sigma,\tau,\sigma\tau=\tau\sigma\}$,其中
\[
  I(\sqrt{-1})=\sqrt{-1},\quad I(\sqrt{3})=\sqrt{3},\quad I(\sqrt{-3})=\sqrt{-3},
\]
\[
  \sigma\tau(\sqrt{-1})=-\sqrt{-1},\quad \sigma\tau(\sqrt{3})=-\sqrt{3},\quad \sigma\tau(\sqrt{-3})=\sqrt{-3}.
\]
由于 $G$ 是 Abel 群,从而 $K/\Q$ 为 Abel 扩张,$G$ 有 $3$ 个 $2$ 阶子群 $\langle\sigma\rangle$,$\langle\tau\rangle$ 和 $\langle\sigma\tau\rangle$,分别对应于 $K$ 的三个中间域 $\Q(\sqrt{3})$,$\Q(\sqrt{-1})$ 和 $\Q(\sqrt{-3})$.记
\[
  K_1=\Q(\sqrt{3}),\qquad K_2=\Q(\sqrt{-3}),\qquad K_3=\Q(\sqrt{-1}).
\]
由定理 \ref{thm:2.8} 知素数 $2$ 在 $K_1$ 中分歧,在 $K_2$ 中不分歧(因 $d(K_2)=-3$).由引理 \ref{lem:ch2-17} 可知 $2$ 在 $K$ 中的惯性域为 $K_2$.又由定理 \ref{thm:2.8} 知 $2$ 在 $K_2$ 中惯性(因 $-3\equiv 5\ (\mathrm{mod}\ 8)$),从而 $2$ 在 $K$ 中的分解域为 $\Q$.类似地,素数 $3$ 在 $K_1$ 中分歧而在 $K_3$ 中不分歧(因 $d(K_3)=-4$),从而 $3$ 的惯性域(即最大不分歧中间域)为 $K_3$.再由 $\left(\dfrac{-1}{3}\right)=-1$ 可知 $3$ 的分解域为 $\Q$.

以下设素数 $p\geqslant 5$.$p$ 在 $K_1$ 和 $K_2$ 中均不分歧,从而在 $K$ 中也不分歧.于是 $p$ 的惯性域为 $K$.如果 $\left(\dfrac{-1}{p}\right)=\left(\dfrac{3}{p}\right)=1$(由二次互反律可知这相当于 $p\equiv 1\ (\mathrm{mod}\ 12)$),则 $p$ 在 $K_1$ 和 $K_3$ 中均分裂,从而在 $K=K_1K_3$ 中完全分裂,即 $p$ 的分解域为 $K$.如果 $\left(\dfrac{-1}{p}\right)=-1$ 而 $\left(\dfrac{3}{p}\right)=1$(即 $p\equiv 11\ (\mathrm{mod}\ 12)$),则 $p$ 在 $K_3$ 中惯性而在 $K_1$ 中分裂,从而 $p$ 的分解域(即使得 $p$ 完全分裂的最大中间域)为 $K_1$.类似地可知,如果 $\left(\dfrac{-1}{p}\right)=1$ 而 $\left(\dfrac{3}{p}\right)=-1$(即 $p\equiv 5\ (\mathrm{mod}\ 12)$),则 $p$ 的分解域为 $K_3$.最后若 $\left(\dfrac{-1}{p}\right)=\left(\dfrac{3}{p}\right)=-1$(即 $p\equiv 7\ (\mathrm{mod}\ 12)$),则 $\left(\dfrac{-3}{p}\right)=1$,而 $p$ 的分解域为 $K_2$.
\end{example}

\subsection{素数在分圆域中的分解}\label{sec:ch2-s3-4}

作为 Frobenius 自同构的又一个应用,现在我们给出分圆域中素理想分解的一个完全刻画.我们已经知道,对于分圆域 $K=\Q(\zeta_m)$,$\zeta_m=\mathrm{e}^{2\pi\mathrm{i}/m}$,$m\not\equiv 2\ (\mathrm{mod}\ 4)$,$K/\Q$ 是 $\varphi(m)$ 次伽罗瓦扩张,并且其伽罗瓦群 $\Gal(K/\Q)$ 自然同构于 $(\Z/m\Z)^\times$:$\sigma_a\mapsto\bar{a}$,$((a,m)=1)$,其中 $\sigma_a$ 是 $\Gal(K/\Q)$ 中满足 $\sigma_a(\zeta_m)=\zeta_m^a$ 的元素,而 $\bar{a}$ 为模 $m$ 剩余类,于是 $K/\Q$ 是 Abel 扩张.$\zeta_m$ 在 $\Q$ 上的极小多项式是分圆多项式
\[
  \Phi_m(x)=\prod_{\substack{1\leqslant r\leqslant m\\ (r,m)=1}}(x-\zeta_m^r).
\]
最后,由伽罗瓦理论易知 $\Q(\zeta_m)\Q(\zeta_n)=\Q(\zeta_{[m,n]})$,$\Q(\zeta_m)\cap\Q(\zeta_n)=\Q(\zeta_{(m,n)})$.

\begin{theorem}\label{thm:2.14}
设 $K=\Q(\zeta_m)$,$m\not\equiv 2\ (\mathrm{mod}\ 4)$,则

(a) 素数 $p$ 在 $K$ 中不分歧的充要条件是 $p\nmid m$.并且在这个时候,
\[
  pO_K=\mathfrak{p}_1\cdots\mathfrak{p}_g,
\]
其中 $g=\varphi(m)/f$,而 $f=f(\mathfrak{p}_i/p)$ 等于 $p\ (\mathrm{mod}\ m)$ 的阶数(即 $f$ 是满足 $p^f\equiv 1\ (\mathrm{mod}\ m)$ 的最小正整数).

(b) 如果 $p\mid m$,令 $m=p^r\cdot m'$,$p\nmid m'$,则
\[
  pO_K=(\mathfrak{p}_1\cdots\mathfrak{p}_g)^e,
\]
其中 $e=\varphi(p^r)$,$g=\varphi(m')/f$,$f=f(\mathfrak{p}_i/p)$ 为 $p\ (\mathrm{mod}\ m')$ 的阶数.
\end{theorem}

\begin{proof}
(a) 第一个推断是由 Dedekind 判别式定理和
\[
  |d(K)|=m^{\varphi(m)}\Big/\prod_{p\mid m}p^{\varphi(m)/(p-1)}
\]
(第一章定理 \ref{thm:1.12})推出,至于第二个推断,我们只需求出 $f(\mathfrak{p}/p)$($\mathfrak{p}=\mathfrak{p}_1$)的值即可.但是 $f(\mathfrak{p}/p)$ 是 Frobenius 自同构的阶,而 $\left(\dfrac{K/\Q}{\mathfrak{p}}\right)$ 由
\[
  \Big(\frac{K/\Q}{\mathfrak{p}}\Big)\alpha\equiv\alpha^p\ (\mathrm{mod}\ \mathfrak{p})\quad (\alpha\in O_K=\Z[\zeta_m])
\]
所刻画.特别地有
\begin{equation}  \Big(\frac{K/\Q}{\mathfrak{p}}\Big)\zeta_m\equiv\zeta_m^p\ (\mathrm{mod}\ \mathfrak{p}).\tag{*}\end{equation}
注意 $\left(\dfrac{K/\Q}{\mathfrak{p}}\right)$ 是 $\Gal(K/\Q)$ 中的元素,从而 $\left(\dfrac{K/\Q}{\mathfrak{p}}\right)\zeta_m=\zeta_m^l$,$(l,m)=1$,因此 $\left(\dfrac{K/\Q}{\mathfrak{p}}\right)$ 的阶即是满足 $l^f\equiv 1\ (\mathrm{mod}\ m)$ 的最小正整数 $f$.为此我们只需再证 $l\equiv p\ (\mathrm{mod}\ m)$ 即可.由 (*) 式我们现在有 $\zeta_m^l\equiv\zeta_m^p\ (\mathrm{mod}\ \mathfrak{p})$.令 $P(x)=x^m-1=\prod_{i=0}^{m-1}(x-\zeta_m^i)$.由于 $p\nmid m$ 而 $O_K/\mathfrak{p}$ 是特征 $p$ 的有限域,从而将 $P(x)$ 和 $P'(x)=mx^{m-1}$ 看作是 $O_K/\mathfrak{p}[x]$ 中多项式时,$P'(x)\neq 0$,并且 $(P(x),P'(x))=1$.这表明 $x^m-1$ 在域 $O_K/\mathfrak{p}$ 上没有重根,换句话说,$\zeta_m$ 看作是 $O_K/\mathfrak{p}$ 中元素时阶也是 $m$.于是由 $\zeta_m^l\equiv\zeta_m^p\ (\mathrm{mod}\ \mathfrak{p})$ 得出 $l\equiv p\ (\mathrm{mod}\ m)$.从而证明了 $f(\mathfrak{p}/p)$(即 $\left(\dfrac{K/\Q}{\mathfrak{p}}\right)$ 的阶)等于满足 $p^f\equiv 1\ (\mathrm{mod}\ m)$ 的最小正整数 $f$.

(b) 令 $K_0=\Q(\zeta_{p^r})$,$K'=\Q(\zeta_{m'})$,则有如下的图示关系:
\[
  \begin{array}{c}
    K=\Q(\zeta_m)\\
    K_0=\Q(\zeta_{p^r})\quad K'=\Q(\zeta_{m'})\\
    \Q
  \end{array}
\]
我们在第 2.4 小节的例 \ref{ex:ch2-s2-4-2} 中已经证明了素数 $p$ 在 $K_0$ 中完全分歧,即 $pO_{K_0}=\mathfrak{p}_0^{\varphi(p^r)}$.从而若 $\mathfrak{p}$ 为 $K$ 中的素理想,$\mathfrak{p}\mid\mathfrak{p}_0$,则 $e(\mathfrak{p}\mid p)\geqslant e(\mathfrak{p}_0/p)=\varphi(p^r)$.另一方面,由 (a) 可知 $p$ 在 $K'$ 中的分解为 $pO_{K'}=\mathfrak{p}_1'\cdots\mathfrak{p}_g'$,其中 $gf=\varphi(m')$,而 $f$ 为 $p$ 模 $m'$ 的阶数.由于
\begin{align*}
  \varphi(m)&=[K:\Q]\geqslant f(\mathfrak{p}\mid p)e(\mathfrak{p}\mid p)g\\
  &\geqslant f\cdot\varphi(p^r)g=\varphi(m')\varphi(p^r)=\varphi(m),
\end{align*}
以及 $f(\mathfrak{p}\mid p)\geqslant f$,$e(\mathfrak{p}\mid p)\geqslant\varphi(p^r)$,可知必然有 $f(\mathfrak{p}/p)=f$,$e(\mathfrak{p}/p)=\varphi(p^r)$.
\end{proof}

\begin{remark}
(1) 我们还可用定理 \ref{thm:2.7} 来证明定理 \ref{thm:2.14} 的 (a).以下记 $\zeta=\zeta_m$,$s=\varphi(m)$,$K=\Q(\zeta)$.分圆多项式 $\Phi_m(x)$(见定理 \ref{thm:1.6})是 $\zeta$ 在 $\Q$ 上的极小多项式,由于 $O_K=\Z[\zeta]$ 有整基 $1,\zeta,\zeta^2,\cdots,\zeta^{s-1}$,在定理 \ref{thm:2.7} 中取 $\alpha=\zeta$,可知每个素数 $p$ 在 $K$ 中素理想分解的参数 $e,f,g$ 均可由 $\Phi_m(x)$ 在 $\F_p[x]$ 中的分解性状所决定.当 $p\nmid m$ 时,$p$ 在 $K$ 中不分歧,于是 $e=1$.而 $\Phi_m(x)$ 在 $\F_p[x]$ 中应当分解成 $g$ 个不同的不可约多项式 $f_i(x)$ 的乘积,并且 $\deg f_i(x)=f$($1\leqslant i\leqslant g$):
\[
  \Phi_m(x)=f_1(x)f_2(x)\cdots f_g(x)\ (\mathrm{mod}\ p),\qquad fg=\varphi(m).
\]
现在取 $\varepsilon$ 为 $\Phi_m(x)$ 在 $\F_p$ 的代数闭包 $\Omega_p$ 中的一个根.上面已证得 $\varepsilon$ 在 $\Omega_p$ 中的阶仍为 $m$.不妨设 $\varepsilon$ 是 $f_1(x)$ 的根,$f_1(x)$ 作为 $\F_p[x]$ 中不可约多项式,在 $\Omega_p$ 中应当有 $f$ 个不同的根 $\varepsilon,\varepsilon^p,\varepsilon^{p^2},\cdots,\varepsilon^{p^{f-1}}$,并且 $f=\deg f_1(x)$ 是满足 $\varepsilon^{p^f}=\varepsilon$ 的最小正整数,由于 $\varepsilon$ 的阶为 $m$,从而 $f$ 是满足 $p^f\equiv 1\ (\mathrm{mod}\ m)$ 的最小正整数.这就证明了定理 \ref{thm:2.14} (a).

(2) 由定理 \ref{thm:2.14} 可知,若 $p$ 为素数,$m=p^rm'$,$(m',p)=1$,则 $p$ 在 $\Q(\zeta_m)$ 中的惯性域为 $\Q(\zeta_{m'})$,而 $p$ 在 $\Q(\zeta_m)$ 中的分解域 $K_D$ 即为它在 $\Q(\zeta_{m'})$ 中的分解域.将 $K_D$ 看作扩张 $\Q(\zeta_{m'})/\Q$ 的中间域,则 $K_D$ 的固定子群是 $\Gal(\Q(\zeta_{m'})/\Q)$ 中由元素
\[
  \Big(\frac{\Q(\zeta_{m'})/\Q}{p}\Big)=\sigma_p
\]
生成的 $f$ 阶循环群 $N=\langle\sigma_p\rangle$,从而 $K_D$ 是 $N$ 的固定子域,$[K_D:\Q]=\varphi(m')/f=g$.
\end{remark}

\begin{example}\label{ex:ch2-s3-4-1}
决定素数 $2$ 在 $K=\Q(\zeta_{28})$ 中的惯性域和分解域.

\emph{解} 由上述知 $2$ 在 $K$ 中的惯性域为 $\Q(\zeta_7)$.以下记 $\zeta=\zeta_7$.而 $2$ 在 $K$ 中的分解域 $K_D$ 即是 $2$ 在 $\Q(\zeta)$ 中的分解域,从而只需考虑扩张 $\Q(\zeta)/\Q$.它的伽罗瓦群为 $6$ 阶循环群
\[
  G=\Gal(\Q(\zeta)/\Q)=\{\sigma_a\mid 1\leqslant a\leqslant 6\}\cong(\Z/7\Z)^\times.
\]
由于 $2\ (\mathrm{mod}\ 7)$ 的阶是 $f=3$,从而 $[K_D:\Q]=g=\frac{6}{3}=2$,即 $K_D=\Q(\sqrt{d})$,其中 $d$ 是某个无平方因子整数.由于 $G$ 是循环群,所以 $K_D$ 也是 $\Q(\zeta)$ 的惟一的二次子域.我们用判别式定理可决定 $d$:
\[
  \begin{array}{ccc}
    \Q(\zeta) & & \{1\}\\
    \Big\vert\ (3) & & \Big\vert\\
    K_D & & N=\langle\sigma_2\rangle\\
    \Big\vert\ (2) & & \Big\vert\\
    \Q & & G\cong(\Z/7\Z)^\times=\langle 3\rangle
  \end{array}
\]
在 $\Q(\zeta)$ 中分歧的素数只有 $7$,而 $7$ 在 $\Q(\zeta)$ 中完全分歧,所以只有 $7$ 在 $K_D$ 中分歧.所以 $\Q(\sqrt{d})$ 的判别式只能为 $\pm 7$,即 $d=\pm 7$.但是 $\Q(\sqrt{7})$ 的判别式为 $28$,因此 $d=-7$,即 $K_D=\Q(\sqrt{-7})$.

我们还可用伽罗瓦理论给出决定 $K_D$ 的另一种方法,它更有一般性.$K_D$ 的固定子群 $N$ 是 $G$ 的 $f=3$ 阶子群,$N$ 是由 $\sigma_2$ 生成的.因此,
\[
  N=\{I=\sigma_1,\sigma_2,\sigma_2^2=\sigma_4\}\quad (\text{而 }\sigma_2^3=\sigma_8=\sigma_1=I).
\]
而 $K_D$ 是 $N$ 的固定子域,即是 $\sigma_2$ 的固定子域.

$\Q(\zeta)$ 中元素 $\alpha=\zeta+\zeta^2+\zeta^4$ 在 $\sigma_2$ 作用下不变:
\[
  \sigma_2(\alpha)=\sigma_2(\zeta)+\sigma_2(\zeta^2)+\sigma_2(\zeta^4)=\zeta^2+\zeta^4+\zeta=\alpha,
\]
可知 $\alpha\in K_D$,即 $\Q(\alpha)\subseteq K_D$.我们来证 $\Q(\alpha)=K_D$.为此只需证明 $\alpha$ 有两个不同的共轭元素即可(因这时 $2=[\Q(\alpha):\Q]\leqslant[K_D:\Q]=2$,从而 $\Q(\alpha)=K_D$).注意
\[
  \sigma_3(\alpha)=\sigma_3(\zeta)+\sigma_3(\zeta^2)+\sigma_3(\zeta^4)=\zeta^3+\zeta^6+\zeta^5=\bar{\alpha}
\]
是 $\alpha$ 的共轭元素.由于 $\zeta,\zeta^2,\cdots,\zeta^6$ 是 $\Z[\zeta]$ 的一组整基,可知 $\alpha=\zeta+\zeta^2+\zeta^4$ 和 $\bar{\alpha}=\zeta^3+\zeta^6+\zeta^5$ 是 $\Z$-线性无关的.特别地,$\alpha\neq\bar{\alpha}$.从而 $\alpha$ 有两个不同的共轭元素 $\alpha$ 和 $\bar{\alpha}$,于是 $K_D=\Q(\alpha)$.

为了看出 $\Q(\alpha)$ 就是 $\Q(\sqrt{-7})$,我们只需算出 $\alpha$ 的值.$\alpha$ 的极小多项式是以 $\alpha$ 和 $\bar{\alpha}$ 为根的二次多项式
\[
  f(x)=(x-\alpha)(x-\bar{\alpha})=x^2-(\alpha+\bar{\alpha})x+\alpha\bar{\alpha}.
\]
由于
\[
  \alpha+\bar{\alpha}=\zeta+\zeta^2+\zeta^3+\zeta^4+\zeta^5+\zeta^6=-1,
\]
\[
  \alpha\bar{\alpha}=(\zeta+\zeta^2+\zeta^4)(\zeta^3+\zeta^5+\zeta^6)=2,
\]
可知 $\alpha,\bar{\alpha}$ 是 $f(x)=x^2+x+2$ 的两个根,即 $\alpha,\bar{\alpha}=\frac{-1\pm\sqrt{-7}}{2}$.所以 $K_D=\Q(-1+\sqrt{-7})=\Q(\sqrt{-7})$.

总起来我们知道:$2$ 在 $K_D=\Q(\sqrt{-7})$ 中分解成两个不同的素理想乘积:$2=\mathfrak{p}_1\mathfrak{p}_2$,而 $\mathfrak{p}_1$ 和 $\mathfrak{p}_2$ 从 $\Q(\sqrt{-7})$ 到 $\Q(\zeta_7)$ 是惯性的,从 $\Q(\zeta_7)$ 到 $\Q(\zeta_{28})=K$ 则完全分歧:$\mathfrak{p}_1=\mathfrak{P}_1^3$,$\mathfrak{p}_2=\mathfrak{P}_2^3$.最后得到
\[
  2O_K=(\mathfrak{P}_1\mathfrak{P}_2)^3,\qquad \mathfrak{P}_1\neq\mathfrak{P}_2.
\]
\end{example}

% ============ 习题 ============
\begin{Exercises}
  \item 证明第 3.1 节中定义的理想范 $N_{L/K}(\mathfrak{a})$ 满足那里的性质 (I)--(IV).
  \item 设 $K/\Q$ 为数域的 Abel 扩张,$K_I$ 是素数 $p$ 对于扩张 $K/\Q$ 的惯性域.求证 $K_I$ 是 $K$ 中使 $p$ 不分歧的最大子域.换句话说,如果 $M$ 是 $K/\Q$ 的中间域,则 $p$ 在 $M$ 中不分歧 $\Leftrightarrow M\subseteq K_I$.
  \item 设 $K=\Q(\sqrt{5},\sqrt{-1})$.求证:
        (a) $K/\Q$ 是 $4$ 次伽罗瓦扩张.求该扩张的伽罗瓦群;
        (b) 有理素数中只有 $2$ 和 $5$ 在 $K$ 中分歧,并且分歧指数均是 $2$;
        (c) 求 $2$ 和 $5$ 对于扩张 $K/\Q$ 的分解群,分解域,惯性群和惯性域;
        (d) $\Q(\sqrt{-5})\subseteq K$,并且 $\Q(\sqrt{-5})$ 中每个素理想在 $K$ 中均不分歧.
  \item 仿照上题的方法,试证明 $\Q(\sqrt{10})\subseteq\Q(\sqrt{2},\sqrt{5})$,并且 $\Q(\sqrt{10})$ 中每个素理想在域 $\Q(\sqrt{2},\sqrt{5})$ 中均不分歧.
  \item 设 $\alpha$ 是多项式 $x^3-x+1$ 的实根,$\beta,\bar{\beta}$ 是此多项式的一对共轭复根,$K=\Q(\alpha)$.求证:
        (a) $\beta+\bar{\beta}=-\alpha$,$\beta\bar{\beta}=-\alpha^{-1}$,
        \[
          [(\beta-\alpha)(\bar{\beta}-\alpha)(\beta-\bar{\beta})]^2=-23;
        \]
        (b) $K(\sqrt{-23})=\Q(\alpha,\beta,\bar{\beta})$,并且 $K(\sqrt{-23})/\Q$ 是 $6$ 次伽罗瓦扩张.求该扩张的伽罗瓦群;
        (c) 决定 $K(\sqrt{-23})$ 的全部子域.其中哪些是 $\Q$ 的伽罗瓦扩域?
        (d) 有理素数中只有 $23$ 在 $K(\sqrt{-23})$ 中分歧,并且分歧指数为 $2$(提示:$23$ 在 $K$ 中的素理想分解式为 $23O_K=\mathfrak{p}_1\mathfrak{p}_2^2$);
        (e) $K(\sqrt{-23})/\Q(\sqrt{-23})$ 是 $3$ 次循环扩张,并且 $\Q(\sqrt{-23})$ 中每个素理想在 $K(\sqrt{-23})$ 中均不分歧.
\end{Exercises}

\begin{remark}
类域论中的一个著名结果是说:每个数域 $K$ 均有极大 Abel 不分歧扩张 $H_K$,称作是 $K$ 的 \emph{Hilbert 类域}.换句话说,$H_K$ 满足如下的性质:(i) $H_K/K$ 是数域的 Abel 扩张;(ii) $K$ 中每个素理想在 $H_K$ 中均不分歧;(iii) 如果 $L$ 为 $K$ 的另一个 Abel 扩域,并且 $K$ 中每个素理想在 $L$ 中均不分歧,则 $L\subseteq H_K$.类域论中还证明了:扩张 $H_K/K$ 的伽罗瓦群同构于数域 $K$ 的理想类群 $C_K$(关于理想类群 $C_K$ 的定义见第三章).我们在第三章中将要证明:域 $K=\Q(\sqrt{-5})$,$\Q(\sqrt{10})$ 和 $\Q(\sqrt{-23})$ 的理想类群分别是 $2$,$2$,$3$ 阶循环群.于是从习题 3--5 可知 $L=\Q(\sqrt{5},\sqrt{-1})$,$\Q(\sqrt{2},\sqrt{5})$ 和 $\Q(\sqrt{-23},\alpha)$($\alpha$ 是多项式 $x^3-x+1$ 的实根)分别是它们的 Hilbert 类域.在一般情形,求一个数域 $K$ 的 Hilbert 类域是代数数论中一个相当困难的问题.
\end{remark}

\begin{enumerate}[1.]
  \setcounter{enumi}{5}
  \item 令 $K=\Q(\zeta)$,$\zeta=\zeta_5=\mathrm{e}^{2\pi\mathrm{i}/5}$.求证 $d(K)=5^4$,并且求 $p=2,3,5,19$ 在 $K$ 中的素理想分解式.
  \item 设 $K=\Q(\omega)$,$\omega=\zeta_{25}$.
        (a) 求证 $K$ 有惟一的一个 $5$ 次子域 $M$;
        (b) 求 $p=2,3,5$ 在 $M$ 中的素理想分解式和它们的分解域与惯性域;
        (c) 求证:素数 $p$ 在 $M$ 中完全分裂 $\Leftrightarrow p\equiv\pm 1,\pm 7\ (\mathrm{mod}\ 25)$.
\end{enumerate}
